% =====================================================================
% Fix 1 - Causal coherence intervention (NeurIPS revision)
% Addresses W1: causal-vs-correlational concern.
% Pre-registration: experiments/fix_01_log.md
% Script:           experiments/fix_01_causal_matched_pairs.py
% Data:             data/revision/fix_01/
% Results:          results/revision/fix_01/
% =====================================================================

\subsection{Causal Intervention: Matched-Similarity HIGH/LOW CCS}
\label{sec:causal_intervention}

% TRACKED-CHANGES (revision): NEW SUBSECTION for \S4.4.
% Why: the prior evidence showed that CCS predicts faithfulness, but did
% not isolate CCS from per-passage query similarity. The reviewer asked
% whether the mechanism is causal or merely correlational.

The preceding results establish that lower Context Coherence Score
(CCS) co-occurs with lower faithfulness. They do not by themselves
establish that coherence is the active factor: per-passage relevance
could be a hidden confound. We therefore add a controlled intervention
that holds mean query similarity fixed while varying CCS.

\paragraph{Construction.}
For each \textsc{SQuAD} query, we retrieve the top-$20$ passages with
the same MiniLM/Chroma index used in the main experiments. We enumerate
all $\binom{20}{3}=1140$ three-passage triples $T$ and compute
\[
  \bar{s}(T)=\frac{1}{3}\sum_{p_i\in T}\cos(q,p_i),
  \qquad
  \mathrm{CCS}(T)=\mu\!\left(\cos(p_i,p_j)\right)
  -\sigma\!\left(\cos(p_i,p_j)\right).
\]
Triples are bucketed into $0.02$-wide bins by $\bar{s}(T)$. Within each
bucket, we choose a HIGH-CCS triple and a LOW-CCS triple that satisfy
$|\bar{s}(T_H)-\bar{s}(T_L)|\le 0.02$,
$\mathrm{CCS}(T_H)-\mathrm{CCS}(T_L)\ge 0.05$, and at most one shared
passage. Among all valid bucket pairs for a query, we keep the pair with
the largest CCS gap. Queries with no valid pair are skipped and logged.

\paragraph{Pre-registered test.}
The primary cell is \textsc{SQuAD}/Mistral-7B with $n\ge 200$ matched
query pairs, seed $42$, and temperature $0.0$. Faithfulness is scored
with the same DeBERTa-v3 NLI scorer used in Table~\ref{tab:hcpc_main}.
The directional test is paired Wilcoxon signed-rank on
$\Delta_i=\mathrm{faith}(T_H^i)-\mathrm{faith}(T_L^i)$ with
\texttt{alternative=greater}. We report Cohen's $d_z$, a
10{,}000-resample bootstrap 95\% CI for the mean paired difference, and
a matched-pair odds ratio for hallucination. The similarity-matching
criterion is enforced directly by construction; a two-sided Wilcoxon
test on similarity deltas is reported only as a diagnostic, not as an
equivalence proof.

\paragraph{Result.}
% TRACKED-CHANGES (revision): FILLED NULL RESULT.
% Why: the preregistered intervention completed with 200 matched pairs
% and did not support the causal HIGH-CCS > LOW-CCS hypothesis. This
% requires downgrading causal/mechanistic language elsewhere.
\begin{table}[t]
\centering
\caption{Fix~1 causal intervention. HIGH-CCS and LOW-CCS triples are
matched within $\pm 0.02$ mean query similarity. The causal claim is
supported only if the paired Wilcoxon $p<0.05$, Cohen's $d_z>0.2$, and
the bootstrap CI for $\bar{\Delta}$ excludes zero.}
\label{tab:fix01_causal}
\footnotesize
\begin{tabular}{lcccccc}
\toprule
Dataset & $n$ & max $|\Delta \bar{s}|$ & CCS gap & $\bar{\Delta}_{faith}$ & $d_z$ & $p$ \\
\midrule
\textsc{SQuAD} & 200 & 0.0185 & 0.5326 & -0.0024 & -0.0171 & 0.6283 \\
\bottomrule
\end{tabular}
\end{table}

The intervention satisfies the matching constraint but does not support
the preregistered causal hypothesis. HIGH-CCS triples have mean
faithfulness $0.6362$, while LOW-CCS triples have mean faithfulness
$0.6386$; the mean paired difference is $-0.0024$ with a bootstrap 95\%
CI of $[-0.0217, 0.0168]$. The one-sided paired Wilcoxon test for
HIGH-CCS $>$ LOW-CCS gives $p=0.6283$ and Cohen's $d_z=-0.0171$.
Hallucination rates are also not directionally favorable to the causal
claim: $0.165$ for HIGH-CCS and $0.090$ for LOW-CCS.

% TRACKED-CHANGES (revision): HONEST DOWNGRADE.
% Why: this null result directly answers the reviewer concern. We keep
% CCS as a diagnostic/predictive signal, but no longer present this
% experiment as evidence that coherence causally drives faithfulness.
We therefore do not interpret CCS as an isolated causal mediator in this
setting. The remaining revision frames CCS as a diagnostic and
deployment-policy signal, and reframes the information-theoretic section
as a consistency check rather than a mechanism proof.
